\documentclass[conference]{IEEEtran}
\IEEEoverridecommandlockouts

\usepackage{cite}
\usepackage{amsmath,amssymb,amsfonts}
\usepackage{graphicx}
\usepackage{textcomp}
\usepackage{xcolor}
\usepackage{booktabs}
\usepackage{array}
\usepackage{url}
\usepackage{hyperref}
\hypersetup{
    colorlinks=true,
    linkcolor=blue,
    citecolor=blue,
    urlcolor=blue
}

\def\BibTeX{{\rm B\kern-.05em{\sc i\kern-.025em b}\kern-.08em
    T\kern-.1667em\lower.7ex\hbox{E}\kern-.125emX}}

\begin{document}

\title{Revisiting Real-World Tool-Use Evaluation for LLM Agents:\\
Gaps, Newer Models, and the Argument Prediction Bottleneck}

\author{
    \IEEEauthorblockN{Muntaha Asim}
    \IEEEauthorblockA{muntahaasim79@gmail.com}
    \and
    \IEEEauthorblockN{Ammar Kashif}
    \IEEEauthorblockA{ammarkashif00@gmail.com}
}

\maketitle

\begin{abstract}
Large language model (LLM) based agents that interact with external tools have emerged as a promising approach toward building general-purpose AI assistants. Evaluating their real-world problem-solving capabilities remains challenging, as existing benchmarks predominantly rely on AI-generated queries, single-step tasks, or text-only interactions that fail to capture the complexity of authentic agent deployments. The GTA benchmark addressed this gap by introducing 229 human-designed, multimodal, real-world queries with executable tool chains spanning perception, operation, logic, and creativity categories. While GTA revealed that even GPT-4 solves fewer than 50\% of tasks---with argument prediction accuracy (ArgAcc) identified as the primary bottleneck---several critical dimensions remain unexplored: the performance of recently released frontier models such as Claude-3.5, Gemini 1.5, DeepSeek-V3, and Llama-3.1, the impact of alternative prompting strategies beyond ReAct, and the role of agent framework selection on end-to-end performance. This paper surveys the existing landscape of LLM tool-use evaluation, systematically identifies these research gaps, and motivates a comprehensive follow-on study targeting the argument prediction bottleneck through structured prompt engineering and agent framework comparison. Our analysis establishes a clear research agenda for advancing the state of the art in real-world LLM tool agents.
\end{abstract}

\begin{IEEEkeywords}
large language models, tool-use agents, benchmark evaluation, argument prediction, prompt engineering, ReAct, agent frameworks
\end{IEEEkeywords}

%=============================================================
\section{Introduction}
%=============================================================

The integration of large language models (LLMs) with external tools has become a central paradigm in the development of capable AI assistants. Frameworks such as LangChain \cite{chase2022langchain}, AutoGPT \cite{gravitas2023autogpt}, and ChatGPT Plugins have demonstrated that equipping LLMs with the ability to invoke specialized tools---ranging from web search and calculators to image processing utilities---dramatically extends the range of tasks these systems can address. At the architectural level, such agent systems decompose complex objectives into two interlocking phases: a \emph{planning} phase, in which the LLM reasons about what actions are required, and an \emph{execution} phase, in which appropriate tools are invoked in sequence. The ReAct framework \cite{yao2023react} formalizes this interplay by interleaving reasoning traces with action steps, and has since become the de facto prompting standard for tool-augmented agents.

Despite this rapid progress, evaluating LLM tool-use in realistic settings remains a fundamentally difficult problem. The dominant benchmarks in this space---ToolBench \cite{qin2023toolbench}, APICBench \cite{patil2023gorilla}, and APIBank \cite{li2023apibank}---each suffer from limitations that undermine their ecological validity. ToolBench relies on ChatGPT to synthesize tool-use queries rather than sourcing them from genuine human needs, resulting in queries that explicitly enumerate required tools and thereby trivialize the planning challenge. APICBench evaluates single-step API matching and does not require multi-tool coordination or execution. APIBank provides a planning evaluation framework but employs virtual rather than real deployed tools, preventing end-to-end execution assessment. Collectively, these benchmarks assess isolated capabilities under artificially simplified conditions that do not reflect the complexity facing a deployed agent.

The GTA benchmark \cite{wang2024gta}, presented at NeurIPS 2024, represents a significant advance by combining three properties absent from prior work: (i) \emph{real user queries}---229 human-written tasks with implicit tool-use steps that require the agent to reason about which tools to invoke; (ii) \emph{real deployed tools}---an executable evaluation platform with 14 tools across perception, operation, logic, and creativity categories, enabling actual end-to-end execution; and (iii) \emph{real multimodal inputs}---authentic images including spatial scenes, web page screenshots, tables, and handwritten materials provided as contextual inputs. Evaluating 16 mainstream LLMs on GTA, the authors find that GPT-4 correctly resolves fewer than 50\% of tasks, while most models fall below 25\%. Fine-grained analysis reveals that argument prediction accuracy (ArgAcc)---the ability to generate correctly formatted tool arguments---is the binding bottleneck, exhibiting the highest Pearson correlation with final answer accuracy among all evaluated metrics.

However, GTA's experimental scope leaves several critical questions unaddressed. First, all evaluations employed a single prompting strategy (ReAct) and a single agent framework (Lagent), leaving open whether architectural or prompt-level changes could alleviate the observed bottlenecks. Second, the tested models represent a snapshot of the field as of early 2024; frontier models released subsequently---including Claude-3.5 Sonnet, Gemini 1.5 Pro, DeepSeek-V3, and Llama-3.1-405B---remain unevaluated on GTA. Third, while AgentFLAN \cite{chen2024agentflan} demonstrates that fine-tuning on ReAct and JSON format data can improve instruction-following and tool-selection accuracy, ArgAcc remains persistently low even after fine-tuning, suggesting that the bottleneck may require targeted prompt-level schema guidance rather than---or in addition to---parameter updates. These gaps collectively motivate the present work.

This paper makes the following contributions. We survey and systematically compare the existing landscape of LLM tool-use benchmarks, characterizing each along five key dimensions and identifying the cumulative evolution that led to GTA. We analyze the evidence for the argument prediction bottleneck and critically assess why existing remedies have proven insufficient. We identify the specific evaluation gaps left open by GTA and articulate a concrete research agenda for extending its evaluation to newer models, alternative prompting strategies, and multiple agent frameworks. Our goal is to provide a rigorous conceptual foundation and actionable guidance for researchers seeking to build more capable and reliable LLM tool agents.

%=============================================================
\section{Related Work}
%=============================================================

\subsection{LLM-Based Agent Systems}

The idea of augmenting LLMs with tool-use capabilities emerged alongside the broader interest in autonomous agents. LangChain \cite{chase2022langchain} established a general-purpose framework for composing LLM calls with tool invocations, enabling developers to construct multi-step pipelines around a wide variety of APIs and data sources. AutoGPT \cite{gravitas2023autogpt} extended this paradigm toward fully autonomous operation, incorporating a persistent planning loop in which the model sets, pursues, and revises goals without per-step human intervention. BabyAGI \cite{nakajima2023babyagi} contributed a task management architecture wherein a controller LLM continuously generates, prioritizes, and executes tasks, demonstrating that even simple prompting loops can produce surprisingly complex emergent behavior.

A pivotal methodological contribution came from Yao et al., who proposed ReAct \cite{yao2023react}, a prompting schema that interleaves natural language \emph{reasoning} traces (``Thought'') with concrete \emph{action} invocations and their observed results. ReAct demonstrated that explicit reasoning traces reduce hallucination, improve interpretability, and enhance task completion relative to action-only or chain-of-thought-only baselines. It quickly became the standard prompting interface for tool-use evaluation. HuggingGPT \cite{shen2024hugginggpt} pushed the paradigm further by using ChatGPT as a central orchestrator that dispatches subtasks to specialized models hosted on HuggingFace, effectively treating the entire HuggingFace model hub as a tool library. While these systems demonstrated the viability of planning-plus-execution decomposition, they were each designed for specific domains or deployment environments. Evaluating their performance in a unified, real-world setting---where tools are genuinely executable and queries are sourced from authentic human needs---remained an open challenge that benchmarking efforts have subsequently sought to address.

\subsection{Tool-Use Benchmarks for LLMs}

The benchmarking of LLM tool-use capabilities has evolved substantially over the past several years. Table~\ref{tab:benchmarks} summarizes the key dimensions along which the major benchmarks differ.

\begin{table}[t]
\caption{Comparison of Tool-Use Benchmarks for LLM Agents}
\label{tab:benchmarks}
\renewcommand{\arraystretch}{1.25}
\resizebox{\columnwidth}{!}{%
\begin{tabular}{lccccc}
\toprule
\textbf{Benchmark} & \textbf{\shortstack{Real\\Queries}} & \textbf{\shortstack{Real\\Tools}} & \textbf{\shortstack{Multi-\\modal}} & \textbf{\shortstack{Human\\Chains}} & \textbf{\shortstack{Exec.\\Eval}} \\
\midrule
APICBench \cite{patil2023gorilla}   & \texttimes & \texttimes & \texttimes & \texttimes & \texttimes \\
ToolBench \cite{qin2023toolbench}   & \texttimes & \checkmark & \texttimes & \texttimes & \texttimes \\
APIBank \cite{li2023apibank}         & \texttimes & \texttimes & \texttimes & \checkmark & \texttimes \\
GAIA \cite{mialon2023gaia}           & \checkmark & \texttimes & \checkmark & \texttimes & \texttimes \\
m\&m's \cite{ma2024mms}             & \texttimes & \checkmark & \checkmark & \checkmark & \texttimes \\
\textbf{GTA} \cite{wang2024gta}      & \checkmark & \checkmark & \checkmark & \checkmark & \checkmark \\
\bottomrule
\end{tabular}%
}
\end{table}

APICBench \cite{patil2023gorilla} introduced the Gorilla LLM, which fine-tunes models to generate accurate API calls drawn from a large API documentation corpus. While it provides broad API coverage, its evaluation tasks are AI-generated and restricted to single-step API selection, precluding assessment of multi-tool coordination. ToolBench \cite{qin2023toolbench} scaled this approach to over 16,000 real-world RESTful APIs and introduced Pass Rate and Win Rate metrics; however, its queries are synthesized by ChatGPT and explicitly mention which tools should be used, eliminating the planning challenge that characterizes authentic tool-use scenarios. APIBank \cite{li2023apibank} addressed the planning dimension by including 53 carefully curated APIs and human-annotated tool chains, but employs virtual rather than real deployed tools, preventing end-to-end execution evaluation.

GAIA \cite{mialon2023gaia} represents a distinct design philosophy: rather than constructing tool-specific tasks, it presents human-written, conceptually simple questions that nonetheless require multi-step reasoning and tool use to answer correctly. GAIA is multimodal and uses real user queries, placing it closer to authentic deployment conditions than the preceding benchmarks. However, it targets general AI assistants rather than tool agents specifically, and does not provide executable tool chains or deployment infrastructure for evaluation. The m\&m's benchmark \cite{ma2024mms} specifically addresses multi-step multimodal planning and provides human-annotated tool chains for evaluation; however, its queries remain AI-generated, limiting its validity as a measure of real-world user intent.

GTA \cite{wang2024gta} is the first benchmark to simultaneously satisfy all five criteria in Table~\ref{tab:benchmarks}. Its 229 tasks span four tool categories---perception, operation, logic, and creativity---and are accompanied by executable tool chains annotated by human experts, enabling fine-grained evaluation at each step of the invocation process. Crucially, GTA evaluates actual end-to-end execution rather than predicted outputs, providing a ground-truth signal that is immune to the post-hoc rationalization artifacts that afflict step-prediction-only evaluation. Nevertheless, as discussed in the preceding section, GTA leaves open the questions of newer model performance, prompting strategy variation, and agent framework comparison.

\subsection{Prompting Strategies for Tool Agents}

Prompting strategy has emerged as a critical design variable in tool-augmented LLM systems. Chain-of-Thought (CoT) prompting \cite{wei2022cot} demonstrated that instructing models to produce intermediate reasoning steps substantially improves performance on multi-step arithmetic, symbolic reasoning, and commonsense tasks. This insight was extended to the tool-use setting by ReAct \cite{yao2023react}, which interleaves CoT-style reasoning with explicit tool invocation actions and observation feedback. Under the ReAct schema, the model emits a ``Thought'' explaining its reasoning, an ``Action'' specifying a tool call and its arguments in a structured format (typically JSON), and observes the returned result before proceeding to the next step. GTA adopts ReAct as its sole evaluation prompting strategy.

GTA's error analysis reveals, however, that the dominant failure mode is not reasoning quality but argument format compliance. As reported in \cite{wang2024gta}, approximately 45.4\% of Llama-3 errors arise from argument predictions that do not adhere to the required JSON schema, while Claude-3's errors are 82.86\% argument-format related despite the model following the overall ReAct structure correctly. This suggests that the bottleneck is not the interleaving of reasoning and action---which ReAct addresses---but the precision of structured argument generation within that schema.

AgentFLAN \cite{chen2024agentflan} attempted to address this through supervised fine-tuning on ReAct-formatted data augmented with JSON-structured tool calls. Results show meaningful improvements in instruction-following accuracy (InstAcc) and tool selection accuracy (ToolAcc), confirming that training signal aligned with the target format helps. However, ArgAcc remains substantially lower than the other metrics even for the fine-tuned Agent-Flan-7B model, indicating that parameter-level updates alone are insufficient to close the argument prediction gap. This points toward the need for prompt-level interventions---such as explicit JSON schema injection, few-shot argument exemplars, or structured output constraints---as complementary or alternative approaches. Evaluating these strategies systematically on GTA constitutes the primary experimental objective of the proposed follow-on study.

\subsection{Multimodal Tool Agents}

The integration of visual modalities into tool-augmented LLM systems introduces an additional layer of complexity that early text-only benchmarks could not assess. Visual ChatGPT \cite{wu2023visualchatgpt} pioneered the combination of visual foundation models with LLMs by routing image-related subtasks to specialized vision models while using the LLM as an orchestrator. MM-ReAct \cite{yang2023mmreact} extended the ReAct prompting paradigm to multimodal settings, enabling ChatGPT to reason over images by composing perception tool calls with language reasoning steps. LLaVA-Plus \cite{liu2023llavaplus} took a learning-based approach, fine-tuning a multimodal LLM to emit tool-use instructions as part of its visual reasoning process. MMLMtool \cite{wang2024mmlmtool} further unified tool learning across modalities within a single generative framework.

These systems collectively establish that multimodal tool use is both practically important and technically non-trivial. GTA's inclusion of authentic image files---spatial scenes, web page screenshots, tables, and handwritten materials---as mandatory contextual inputs represents a more realistic evaluation of multimodal tool-use capability than prior work, which typically uses synthetic or curated image sets. The proposed study preserves this multimodal dimension across all experimental conditions, ensuring that findings remain relevant to the practical deployment scenarios that motivate the research.

%=============================================================
\section{Research Gaps and Motivation}
%=============================================================

The foregoing review surfaces three concrete gaps in the current literature that motivate the experimental program proposed for the subsequent sections of this paper.

\textbf{Gap 1: Unevaluated frontier models.} GTA's model evaluations were conducted in early 2024 and cover models up to GPT-4 and Claude-3. Frontier models released since then---including Claude-3.5 Sonnet, Gemini 1.5 Pro, DeepSeek-V3, and Llama-3.1-405B---differ substantially in their instruction-following capabilities and JSON generation reliability. Extending GTA evaluation to these models is necessary to establish whether the argument prediction bottleneck persists under current state-of-the-art conditions.

\textbf{Gap 2: Single prompting strategy.} GTA uses ReAct exclusively as its prompting interface. Given that the dominant failure mode is argument format error, prompt-level interventions that explicitly supply JSON schema definitions, provide few-shot argument examples, or impose structured output constraints warrant systematic evaluation. Whether such strategies can close the ArgAcc gap without parameter updates is an open empirical question with significant practical implications.

\textbf{Gap 3: Single agent framework.} GTA uses the Lagent framework exclusively for end-to-end evaluation. Agent frameworks differ in how they handle tool registration, error recovery, and multi-turn dialogue management---all of which may interact with the argument prediction bottleneck. A comparative evaluation across frameworks such as LangChain \cite{chase2022langchain}, AutoGPT \cite{gravitas2023autogpt}, and Lagent would establish whether framework choice constitutes a meaningful design variable for practitioners.

Addressing these three gaps through controlled experimentation on the GTA benchmark will yield actionable guidance for researchers and engineers developing real-world LLM tool agents.

%=============================================================
\section{Conclusion}
%=============================================================

This paper has surveyed the landscape of LLM tool-use evaluation, traced the evolution from early single-step API benchmarks to the comprehensive real-world evaluation framework introduced by GTA, and systematically identified three research gaps that the current literature leaves unaddressed. The argument prediction bottleneck---wherein correct tool selection and instruction following are undermined by argument format errors---emerges as the central challenge for advancing LLM tool agents beyond current performance levels. Future work will address this bottleneck through controlled experiments comparing prompting strategies, evaluating newly released frontier models, and assessing the impact of agent framework choice on end-to-end task performance.

%=============================================================
\begin{thebibliography}{00}

\bibitem{qin2023toolbench}
Y. Qin, S. Liang, Y. Ye, K. Zhu, L. Yan, Y. Lu, Y. Lin, X. Cong, X. Tang, B. Qian, S. Zhao, R. Tian, R. Xie, J. Zhou, M. Gerstein, D. Li, Z. Liu, and M. Sun, ``ToolLLM: Facilitating large language models to master 16000+ real-world APIs,'' \textit{arXiv preprint arXiv:2307.16789}, 2023.

\bibitem{li2023apibank}
M. Li, F. Song, B. Yu, H. Yu, Z. Li, F. Huang, and Y. Li, ``API-Bank: A comprehensive benchmark for tool-augmented LLMs,'' in \textit{Proc. Conf. Empirical Methods Natural Language Process. (EMNLP)}, pp.~3102--3116, 2023.

\bibitem{patil2023gorilla}
S.~G. Patil, T. Zhang, X. Wang, and J.~E. Gonzalez, ``Gorilla: Large language model connected with massive APIs,'' \textit{arXiv preprint arXiv:2305.15334}, 2023.

\bibitem{mialon2023gaia}
G. Mialon, C. Fourrier, C. Swift, T. Wolf, Y. LeCun, and T. Scialom, ``GAIA: A benchmark for general AI assistants,'' in \textit{Proc. Int. Conf. Learning Representations (ICLR)}, 2024.

\bibitem{ma2024mms}
Z. Ma, J. He, J. Li, L. Chen, Q. Guo, Y. Wang, and A. Hauptmann, ``m\&m's: A benchmark to evaluate tool use for multi-step multi-modal tasks,'' in \textit{Synthetic Data for Computer Vision Workshop @ CVPR}, 2024.

\bibitem{yao2023react}
S. Yao, J. Zhao, D. Yu, N. Du, I. Shafran, K. Narasimhan, and Y. Cao, ``ReAct: Synergizing reasoning and acting in language models,'' in \textit{Proc. Int. Conf. Learning Representations (ICLR)}, 2023.

\bibitem{chen2024agentflan}
Z. Chen, S. Liu, D. Zhou, R. Feng, X. Wan, and W. Zhang, ``Agent-FLAN: Designing data and methods for effective agent tuning,'' in \textit{Findings Assoc. Comput. Linguistics (ACL)}, pp.~5553--5571, 2024.

\bibitem{wang2024gta}
J. Wang, Z. Ma, Y. Li, S. Zhang, C. Chen, K. Chen, and X. Le, ``GTA: A benchmark for general tool agents,'' in \textit{Proc. 38th Conf. Neural Inf. Process. Syst. (NeurIPS), Datasets and Benchmarks Track}, 2024.

\bibitem{chase2022langchain}
H. Chase, ``LangChain,'' \textit{GitHub Repository}, Oct. 2022. [Online]. Available: \url{https://github.com/langchain-ai/langchain}

\bibitem{gravitas2023autogpt}
S. Gravitas, ``AutoGPT,'' \textit{GitHub Repository}, 2023. [Online]. Available: \url{https://github.com/Significant-Gravitas/AutoGPT}

\bibitem{nakajima2023babyagi}
Y. Nakajima, ``BabyAGI,'' \textit{GitHub Repository}, 2023. [Online]. Available: \url{https://github.com/yoheinakajima/babyagi}

\bibitem{shen2024hugginggpt}
Y. Shen, S. Song, X. Tan, D. Li, W. Lu, and Y. Zhuang, ``HuggingGPT: Solving AI tasks with ChatGPT and its friends in HuggingFace,'' in \textit{Proc. 37th Conf. Neural Inf. Process. Syst. (NeurIPS)}, 2024.

\bibitem{wei2022cot}
J. Wei, X. Wang, D. Schuurmans, M. Bosma, F. Xia, E. Chi, Q.~V. Le, and D. Zhou, ``Chain-of-thought prompting elicits reasoning in large language models,'' in \textit{Proc. 36th Conf. Neural Inf. Process. Syst. (NeurIPS)}, vol.~35, pp.~24824--24837, 2022.

\bibitem{wu2023visualchatgpt}
C. Wu, S. Yin, W. Qi, X. Wang, Z. Tang, and N. Duan, ``Visual ChatGPT: Talking, drawing and editing with visual foundation models,'' \textit{arXiv preprint arXiv:2303.04671}, 2023.

\bibitem{yang2023mmreact}
Z. Yang, L. Li, J. Wang, K. Lin, E. Azarnasab, F. Ahmed, Z. Liu, C. Liu, M. Zeng, and L. Wang, ``MM-ReAct: Prompting ChatGPT for multimodal reasoning and action,'' \textit{arXiv preprint arXiv:2303.11381}, 2023.

\bibitem{liu2023llavaplus}
H. Liu, C. Li, Y. Li, B. Li, Y. Zhang, S. Shen, and Y.~J. Lee, ``LLaVA-Plus: Learning to use tools for creating multimodal agents,'' \textit{arXiv preprint arXiv:2311.05437}, 2023.

\bibitem{wang2024mmlmtool}
J. Wang, W. Lan, Z. Qian, S. Li, X. Hua, T. Su, and M. Luo, ``MMLMtool: A multimodal large language model for tool learning,'' \textit{arXiv preprint arXiv:2401.10727}, 2024.

\end{thebibliography}

\end{document}
